\documentclass[12pt,a4paper]{article}

\usepackage{fontawesome5}
\providecommand{\faPencilAlt}{\faIcon{pencil-alt}}

\usepackage{styles/couleurs}
\usepackage{styles/mise_en_page}
\usepackage{styles/boites}
\usepackage{styles/maths}
\usepackage{styles/figures}

\usepackage{enumitem}
\usepackage{array}
\usepackage{tabularx}
\usepackage{amsmath}
\usepackage{amssymb}
\usepackage{graphicx}
\usepackage{tikz}
\usepackage{xparse}

\usetikzlibrary{calc,arrows.meta}

\setlist[itemize]{leftmargin=1.4em,itemsep=0.15em,topsep=0.2em}
\setlist[enumerate]{leftmargin=*,itemsep=0.22em,topsep=0.25em}

\newcommand{\lignepointillee}{%
  \noindent\dotfill\par
  \vspace{0.45em}
}

\newcommand{\lignesreponse}[1]{%
  \begin{tcolorbox}[
    enhanced,
    colback=white,
    colframe=NeutreC!40,
    boxrule=0.5pt,
    arc=1.5mm,
    left=3mm,
    right=3mm,
    top=2mm,
    bottom=1mm,
    before skip=2pt,
    after skip=2pt
  ]
    \foreach \i in {1,...,#1}{\lignepointillee}
  \end{tcolorbox}
}

\newcommand{\imageactivite}[3]{%
  \IfFileExists{#3}{%
    \includegraphics[width=#1,height=#2,keepaspectratio]{#3}%
  }{%
    \fbox{%
      \parbox[c][#2][c]{#1}{%
        \centering\sffamily\small
        Image à intégrer\par
        \vspace{0.3em}
        \texttt{\detokenize{#3}}
      }%
    }%
  }%
}

\newcommand{\titreseq}[1]{%
\begin{center}
  {\sffamily\LARGE\bfseries\color{TitreC}
  Séquence 5 -- Statistiques et proportionnalité}\\[0.2em]
  {\sffamily\large\bfseries\color{SavoirC}#1}
\end{center}
}

\begin{document}

\titreseq{Activité 3 -- Compléter un tableau de proportionnalité}

\section*{Partie A -- Additions et soustractions}

Un marchand vend des crêpes au même prix à l'unité.

\[
\begin{array}{c|ccccccc}
\text{Nombre de crêpes} & 1 & 2 & 3 & 4 & 5 & 6 & 7\\
\hline
\text{Prix en euros} & & & 1{,}8 & & & & 4{,}2
\end{array}
\]

Complète le tableau en utilisant d'abord uniquement des additions et des soustractions.
Indique au moins deux relations utilisées entre les colonnes.

\lignesreponse{5}

\section*{Partie B -- Multiplications et divisions}

Un marchand vend des galettes au même prix à l'unité.

\[
\begin{array}{c|ccccccc}
\text{Nombre de galettes} & 1 & 2 & 3 & 4 & 5 & 6 & 7\\
\hline
\text{Prix en euros} & & & 2{,}1 & 2{,}8 & & &
\end{array}
\]

Complète le tableau en utilisant des multiplications et des divisions.

\lignesreponse{5}

\section*{Partie C -- Retour à l'unité}

Des pommes sont vendues au poids. Pour \(3\) kg, on paie \(6\) €.

\begin{enumerate}
  \item Quel est le prix de \(1\) kg ?
  \item Complète le tableau.
\end{enumerate}

\[
\begin{array}{c|ccccc}
\text{Masse en kg} & 1 & 2{,}5 & 7 & 1{,}8 & 0{,}175\\
\hline
\text{Prix en euros} & & & & &
\end{array}
\]

\lignesreponse{6}

\section*{Partie D -- Choisir une méthode}

Pour chaque partie, nomme la méthode qui t'a semblé la plus efficace :

\[
\text{additivité}
\quad
\text{multiplication d'une colonne}
\quad
\text{coefficient}
\quad
\text{retour à l'unité}.
\]

\lignesreponse{3}

\end{document}
